%% =========================================================
%% Chapter: Examples of PEPS on the torus
%% =========================================================

\chapter{Examples of PEPS on the torus}%
\label{ch:peps_examples}

The simplest PEPS place one four-leg tensor at every site of a square lattice
on the torus. This chapter fixes that setting and computes the states of the
standard examples of \cite[Appendix~A]{Cirac2021Matrix}: the GHZ state, the
cluster state, the CZX model, the AKLT state, the resonating valence bond
state, and the PEPS whose coefficients are the Boltzmann weights of classical
models. The toric code and the quantum double models are treated with the
theory of $G$-injective tensors in Chapter~\ref{ch:peps_quantum_double}.

\section{Torus PEPS of a four-leg tensor}%
\label{sec:peps_torus_site_tensor}

Two-dimensional PEPS on the discrete torus are defined by a four-leg tensor
$A^i_{\alpha\beta\gamma\delta}$ placed at every site of a square lattice, with
virtual indices ordered top, right, down, left
\cite[Appendix~A, ``Two dimensions: PEPS'']{Cirac2021Matrix}. On a torus, the
four legs of site $v=(x,y)$ correspond to the edges from $v$ to $(x,y+1)$
(top), from $v$ to $(x+1,y)$ (right), from $(x,y-1)$ to $v$ (down), and from
$(x-1,y)$ to $v$ (left).

\begin{definition}[Torus PEPS of a four-leg tensor]%
    \label{def:peps_torus_site_tensor}
    \lean{TNLean.PEPS.torusDownEdge, TNLean.PEPS.torusLeftEdge,
        TNLean.PEPS.torusTopLeg, TNLean.PEPS.torusRightLeg,
        TNLean.PEPS.torusDownLeg, TNLean.PEPS.torusLeftLeg,
        TNLean.PEPS.torusLeftEdge_add_one, TNLean.PEPS.torusDownEdge_add_one,
        TNLean.PEPS.torusSiteTensor}
    \leanok
    \uses{def:peps_tensor, def:peps_torus_uniform_bond_dim}
    Let $a^i_{\alpha\beta\gamma\delta}$ be a tensor with physical dimension
    $d$ and bond dimension $D$. Its torus PEPS on the torus of width $w$ and
    height $h$ has bond dimension $D$ on every edge and, at every site $v$,
    the local tensor $a$ with its four virtual indices read on the top,
    right, down and left legs of $v$
    \cite[Appendix~A, ``Two dimensions: PEPS'']{Cirac2021Matrix}. The right
    leg of $v$ is the left leg of $v+(1,0)$, and the top leg of $v$ is the
    down leg of $v+(0,1)$.
\end{definition}

\begin{lemma}[Shift-invariant functions on the torus are constant]%
    \label{lem:peps_torus_shift_const}
    \lean{TNLean.PEPS.torusVertex_apply_eq_apply_zero_of_shift}
    \leanok
    \uses{def:peps_torus_geometry_translation}
    If a function $f$ on the sites of the torus satisfies $f(x+1,y)=f(x,y)$
    and $f(x,y+1)=f(x,y)$ for all $(x,y)$, then $f$ is constant.
\end{lemma}

\begin{proof}\leanok
    Induction on the representatives $0,\ldots,w-1$ of the first coordinate
    and $0,\ldots,h-1$ of the second.
\end{proof}

\begin{lemma}[Contraction of a torus PEPS]%
    \label{lem:peps_torus_site_tensor_contraction}
    \lean{TNLean.PEPS.stateCoeff_torusSiteTensor_eq_sum_edge,
        TNLean.PEPS.isHorizontalTorusEdge_torusRightEdge,
        TNLean.PEPS.isVerticalTorusEdge_torusUpEdge,
        TNLean.PEPS.torusRightEdge_injective, TNLean.PEPS.torusUpEdge_injective,
        TNLean.PEPS.torusRightEdge_ne_torusUpEdge, TNLean.PEPS.torusEdgeEquiv,
        TNLean.PEPS.stateCoeff_torusSiteTensor}
    \leanok
    \uses{def:peps_torus_site_tensor, def:peps_stateCoeff}
    Let $w,h\ge 3$. The right edges and the up edges of the sites are
    pairwise distinct and exhaust the edges of the torus, and the coefficient
    of the torus PEPS of $a$ is
    \begin{align}
        \Coeff(\sigma)=\sum_{h_\bullet,v_\bullet}\ \prod_{(x,y)}
        a^{\sigma(x,y)}_{v_{(x,y)}\,h_{(x,y)}\,v_{(x,y-1)}\,h_{(x-1,y)}},
        \notag
    \end{align}
    the sum running over all assignments of bond indices $h_{(x,y)}$ to the
    horizontal edges from $(x,y)$ to $(x+1,y)$ and $v_{(x,y)}$ to the
    vertical edges from $(x,y)$ to $(x,y+1)$.
\end{lemma}

\begin{proof}\leanok
    Every edge is horizontal or vertical, and every horizontal edge is the
    right edge of one of its endpoints and every vertical edge the up edge
    of one of its endpoints. For $w\ge 3$ the right edges of distinct sites
    are distinct: otherwise $p=q+(1,0)$ and $q=p+(1,0)$, so $2=0$ in the
    horizontal coordinate. The same holds vertically for $h\ge 3$. A
    horizontal edge is never vertical. Reindexing the sum over edge
    labellings along this bijection gives the formula.
\end{proof}

On a torus of width two the right and left neighbours of a site coincide and
the two edges joining them are one edge of the lattice graph, so the four
legs of a site are not four distinct bonds. The contraction formula, and the
cluster, CZX, AKLT, RVB and classical-model statements below, therefore require
$w,h\ge 3$; the
source constructions are local and apply to every torus
\cite{gap:rmp_peps_examples_small_torus}.

A virtual string on a torus PEPS acts on the bonds it crosses, diagonally in
the bond basis: a bond carrying the index $\alpha$ contributes an element
$X_\alpha$ of a monoid $M$, and the contributions are multiplied in the order
in which the string crosses the bonds. For a matrix algebra $M$ this is a
matrix product operator on the virtual level.

\begin{lemma}[Pulling a virtual string through a rectangle]%
    \label{lem:peps_torus_string_pull_through}
    \lean{TNLean.PEPS.gridPath_column_prod_eq, TNLean.PEPS.gridPath_prod_eq,
        TNLean.PEPS.torusWestSouthString, TNLean.PEPS.torusNorthEastString,
        TNLean.PEPS.torusWestSouthString_eq_torusNorthEastString,
        TNLean.PEPS.sum_smul_torusWestSouthString_eq}
    \leanok
    \uses{def:peps_torus_site_tensor}
    Let $a$ be a four-leg tensor and $\alpha\mapsto X_\alpha$ a map from the
    bond indices to a monoid $M$ with $X_tX_r=X_lX_b$ whenever
    $a^{s}_{trbl}\neq0$. Fix a rectangle of $m$ columns and $n$ rows of sites
    of the torus. Consider the string from its upper left to its lower right
    corner that crosses the left legs of the leftmost column, from top to
    bottom, and then the down legs of the bottom row, from left to right, and
    the string that crosses the top legs of the top row and then the right
    legs of the rightmost column. On every labelling of the bonds whose term
    in the contraction of the torus PEPS of $a$ is nonzero the two strings
    have the same product, and so, for a semiring $M$ with a $\C$-module structure, the
    contracted network with the first string inserted equals the contracted
    network with the second string inserted. This is the pulling-through
    equation of a matrix product operator symmetry
    \cite[Section~``MPO and PEPO'']{Cirac2021Matrix} for strings that are
    diagonal in the bond basis.
\end{lemma}

\begin{proof}\leanok
    Write $P_{ij}$ for the contribution of the left leg of the site in
    column $i$ and row $j$ of the rectangle, counted from its upper left
    site, and $Q_{ij}$ for that of its top leg. The right leg of that site is
    the left leg of the site in column $i+1$, and its down leg is the top leg
    of the site in row $j+1$, so the hypothesis at a site with nonzero entry
    reads $Q_{ij}P_{i+1,j}=P_{ij}Q_{i,j+1}$. Induction on $n$ gives
    $P_{i0}\cdots P_{i,n-1}\,Q_{in}=Q_{i0}\,P_{i+1,0}\cdots P_{i+1,n-1}$
    for one column, and induction on $m$, moving one column at a time, gives
    the identity for the rectangle. A nonzero term of the contraction has a
    nonzero entry at every site, and the terms that vanish contribute zero
    to both sides.
\end{proof}

\section{The GHZ state}%
\label{sec:peps_ghz}

The two-dimensional GHZ state is a superposition of the $d$ constant
spin configurations \cite[Appendix~A, ``The GHZ state'']{Cirac2021Matrix}.
All bonds are contracted periodically. Separate horizontal and vertical
assignments of bond indices retain the self-loops and parallel bonds at small lattice sizes.
The simple-graph representation requires $w,h\ge2$
\cite{gap:rmp_peps_examples_small_torus}.

\begin{definition}[Two-dimensional GHZ tensor]%
    \label{def:peps_ghz_tensor}
    \lean{TNLean.PEPS.ghzSiteTensor, TNLean.PEPS.ghzPEPS}
    \leanok
    \uses{def:peps_torus_site_tensor}
    The GHZ tensor has $D=d$ and
    $A^i_{\alpha\beta\gamma\delta}=\delta_{i=\alpha=\beta=\gamma=\delta}$
    \cite[Appendix~A, ``The GHZ state'']{Cirac2021Matrix}.
\end{definition}

\begin{theorem}[The GHZ PEPS is the GHZ state]%
    \label{thm:peps_ghz_state}
    \lean{TNLean.PEPS.sum_prod_ghzSiteTensor, TNLean.PEPS.stateCoeff_ghzPEPS}
    \leanok
    \uses{def:peps_ghz_tensor, def:peps_stateCoeff,
        lem:peps_torus_site_tensor_contraction}
    On every torus with $w,h\ge1$, the GHZ tensor generates the unnormalized
    state $\sum_{i=0}^{d-1}\ket{i,i,\ldots,i}$
    \cite[Appendix~A, ``The GHZ state'']{Cirac2021Matrix}.
    With horizontal and vertical assignments of bond indices $h_{(x,y)},v_{(x,y)}$ as in
    Lemma~\ref{lem:peps_torus_site_tensor_contraction}, its coefficient is
    \begin{align}
        \sum_{h,v}\prod_{(x,y)}
        A^{\sigma(x,y)}_{v_{(x,y)},h_{(x,y)},v_{(x,y-1)},h_{(x-1,y)}}
         & = \sum_i\prod_{(x,y)}\delta_{\sigma(x,y),i}.
        \label{eq:peps_ghz_periodic_contraction}
    \end{align}
    For $w,h\ge2$, the simple-graph GHZ PEPS has the same coefficients.
\end{theorem}

\begin{proof}\leanok
    \uses{lem:peps_torus_shift_const}
    A virtual configuration contributes only if every leg of every site
    carries the physical index of that site. The right leg of $v$ is the left
    leg of $v+(1,0)$ and the top leg of $v$ is the down leg of $v+(0,1)$, so
    $\sigma$ is invariant under both unit shifts and is constant by
    Lemma~\ref{lem:peps_torus_shift_const}. Every edge is then labelled by
    that constant, so exactly one virtual configuration contributes when
    $\sigma$ is constant and none otherwise.
\end{proof}

\begin{theorem}[Non-injectivity of the GHZ PEPS]
    \label{thm:peps_ghz_not_injective}
    \lean{TNLean.PEPS.ghzSiteTensor_not_linearIndependent,
        TNLean.PEPS.ghzPEPS_not_isVertexInjective}
    \leanok
    \uses{def:peps_ghz_tensor, def:IsVertexInjective}
    For $d\ge2$, the four-leg GHZ tensor is not injective.
    For $d\ge2$ and $w,h\ge2$, the simple-graph GHZ PEPS is not vertex-injective.
\end{theorem}

\begin{proof}\leanok
    \uses{lem:peps_torus_site_tensor_contraction}
    The four-leg basis vector with labels $(0,1,0,0)$ has zero image.
    In the simple graph, give the right bond at one vertex label $1$
    and every other incident bond label $0$. The top and right bonds are distinct,
    so no physical label can agree with both. This virtual basis
    vector is mapped to zero.
\end{proof}

\begin{theorem}[On-site permutations and their virtual action]
    \label{thm:peps_ghz_permutation}
    \lean{TNLean.PEPS.ghzSiteTensor_perm, TNLean.PEPS.ghzSiteTensor_add,
        TNLean.PEPS.sum_prod_ghzSiteTensor_perm,
        TNLean.PEPS.stateCoeff_ghzPEPS_perm, TNLean.PEPS.stateCoeff_ghzPEPS_add}
    \leanok
    \uses{def:peps_ghz_tensor, def:peps_stateCoeff, thm:peps_ghz_state}
    Every permutation $\pi$ of the $d$ labels acts on the GHZ site tensor by
    \begin{align}
        A^{\pi(i)}_{\pi(\alpha)\pi(\beta)\pi(\gamma)\pi(\delta)}
         & = A^i_{\alpha\beta\gamma\delta}.
    \end{align}
    For every positive $w,h$, the contraction in
    (\ref{eq:peps_ghz_periodic_contraction}) is unchanged when $\sigma$
    is replaced by $\pi\circ\sigma$. For $w,h\ge2$, the simple-graph
    GHZ PEPS therefore satisfies $\Coeff(\pi\circ\sigma)=\Coeff(\sigma)$.
    In particular, for $d\ge1$ the on-site $\Z_d$ shift is a symmetry,
    and its virtual action is the same shift on every bond.
\end{theorem}

\begin{proof}\leanok
    \uses{thm:peps_ghz_state}
    A permutation preserves equality of the five labels in the site
    tensor. Reindexing the sum over constant configurations in
    Theorem~\ref{thm:peps_ghz_state} proves the state identity.
    Take $\pi(i)=i+g$ for the cyclic shift.
\end{proof}

\section{The cluster state}%
\label{sec:peps_cluster}

\begin{definition}[Two-dimensional cluster tensor]%
    \label{def:peps_cluster_tensor}
    \lean{TNLean.PEPS.clusterSiteTensor, TNLean.PEPS.clusterPEPS}
    \leanok
    \uses{def:peps_torus_site_tensor}
    The cluster tensor is $A=\ket{0}(00{+}{+}|+\ket{1}(11{-}{-}|$ with
    $(\pm|=[(0|\pm(1|]/\sqrt{2}$: the top and right legs copy the physical
    index $s$, and the down and left legs carry
    $(-1)^{sb}/\sqrt{2}$ at bond index $b$
    \cite[Appendix~A, ``The cluster state'']{Cirac2021Matrix}.
\end{definition}

\begin{theorem}[The cluster PEPS]%
    \label{thm:peps_cluster_state}
    \lean{TNLean.PEPS.stateCoeff_clusterPEPS}
    \leanok
    \uses{def:peps_cluster_tensor, def:peps_stateCoeff}
    On a torus with $w,h\ge 3$ and $N=wh$ sites the cluster PEPS has
    coefficient
    \begin{align}
        \Coeff(\sigma)=2^{-N}\prod_{v}(-1)^{\sigma_v\sigma_{v+(1,0)}
                +\sigma_v\sigma_{v+(0,1)}}.
        \notag
    \end{align}
\end{theorem}

\begin{proof}\leanok
    \uses{lem:peps_torus_site_tensor_contraction}
    By Lemma~\ref{lem:peps_torus_site_tensor_contraction} the top and right
    legs force every bond index to equal the physical index of the site
    below or to the left of it. The remaining factor at $v$ is
    $\tfrac12(-1)^{\sigma_v\sigma_{v-(0,1)}+\sigma_v\sigma_{v-(1,0)}}$;
    shifting the product over $v$ by one site in each direction gives the
    formula.
\end{proof}

\begin{definition}[Graph-state circuit and stabilizers]
    \label{def:peps_graph_stabilizers}
    \lean{TNLean.PEPS.graphCircuitSign, TNLean.PEPS.graphFlip,
        TNLean.PEPS.graphNeighborSign, TNLean.PEPS.graphStabilizer,
        TNLean.PEPS.clusterBondTarget, TNLean.PEPS.clusterCircuitState,
        TNLean.PEPS.clusterStabilizer}
    \leanok
    \uses{def:peps_torus_geometry_translation}
    Let a finite set of bonds have distinct endpoints $a(e),b(e)$ in a finite
    vertex set $V$. Parallel bonds are counted with multiplicity. The product
    of controlled-$Z$ gates has diagonal coefficient
    $g(\sigma)=\prod_e(-1)^{\sigma_{a(e)}\sigma_{b(e)}}$.
    The stabilizer $K_v$ flips the qubit at $v$ and applies $Z$ at the opposite
    endpoint of every incident bond. Thus
    \begin{align}
        (K_v\psi)(\sigma)
         & =\prod_{e:a(e)=v}(-1)^{\sigma_{b(e)}}
        \prod_{e:b(e)=v}(-1)^{\sigma_{a(e)}}
        \psi(\sigma\mathbin{\oplus}\mathbf{1}_v).
    \end{align}
    On the square torus, take the horizontal and vertical positive bonds;
    the normalized circuit state has coefficients $2^{-N/2}g(\sigma)$
    \cite[Appendix~A, ``The cluster state'']{Cirac2021Matrix}.
\end{definition}

\begin{theorem}[Common graph-stabilizer eigenspace]
    \label{thm:peps_graph_stabilizer_line}
    \lean{TNLean.PEPS.graphStabilizer_graphCircuitSign,
        TNLean.PEPS.graphStabilizer_fixed_iff}
    \leanok
    \uses{def:peps_graph_stabilizers}
    For the loop-free bond family of Definition~\ref{def:peps_graph_stabilizers},
    $K_v\psi=\psi$ for every $v$ if and only if $\psi=cg$ for some $c\in\C$.
\end{theorem}

\begin{proof}\leanok
    Flipping $v$ multiplies $g$ by the product of the neighbouring $Z$ signs:
    \begin{align}
        g(\sigma\mathbin{\oplus}\mathbf{1}_v)=s_v(\sigma)g(\sigma),
        \qquad
        s_v(\sigma)=\prod_{e:a(e)=v}(-1)^{\sigma_{b(e)}}
        \prod_{e:b(e)=v}(-1)^{\sigma_{a(e)}}.
    \end{align}
    Since $s_v(\sigma)^2=1$, this gives $K_vg=g$. Conversely,
    $K_v\psi=\psi$ reads $\psi(\sigma)=s_v(\sigma)\psi(\sigma\mathbin{\oplus}\mathbf{1}_v)$,
    so
    \begin{align}
        g(\sigma\mathbin{\oplus}\mathbf{1}_v)\psi(\sigma\mathbin{\oplus}\mathbf{1}_v)
        =g(\sigma)\psi(\sigma)
    \end{align}
    for every site $v$. Changing the sites one at a time reduces every
    configuration to the all-zero configuration. Hence
    $g(\sigma)\psi(\sigma)=\psi(0)$ and $\psi(\sigma)=\psi(0)g(\sigma)$.
\end{proof}

\begin{theorem}[Square-torus cluster circuit and uniqueness]
    \label{thm:peps_cluster_stabilizer_line}
    \lean{TNLean.PEPS.stateCoeff_clusterPEPS_eq_smul_graphCircuitSign,
        TNLean.PEPS.stateCoeff_clusterPEPS_eq_smul_clusterCircuitState,
        TNLean.PEPS.clusterStabilizer_apply,
        TNLean.PEPS.clusterStabilizer_fixed_iff,
        TNLean.PEPS.stateCoeff_clusterPEPS_ne_zero}
    \leanok
    \uses{def:peps_cluster_tensor, def:peps_graph_stabilizers}
    On a torus with $w,h\ge3$ and $N=wh$, the printed cluster PEPS is
    $2^{-N/2}$ times the normalized controlled-$Z$ circuit state and is nonzero.
    Its span is precisely the common $+1$ eigenspace of
    \begin{align}
        K_v=X_v Z_{v+(1,0)}Z_{v-(1,0)}Z_{v+(0,1)}Z_{v-(0,1)}.
        \label{eq:peps_cluster_stabilizer_pauli}
    \end{align}
\end{theorem}

\begin{proof}\leanok
    \uses{thm:peps_cluster_state, thm:peps_graph_stabilizer_line}
    Theorem~\ref{thm:peps_cluster_state} identifies the coefficient with
    $2^{-N}g(\sigma)$, and its all-zero coefficient is $2^{-N}\ne0$.
    Each site has one incoming and one outgoing bond in each coordinate
    direction, giving (\ref{eq:peps_cluster_stabilizer_pauli}). Apply
    Theorem~\ref{thm:peps_graph_stabilizer_line} and absorb the nonzero scalar.
\end{proof}

\section{The CZX model}%
\label{sec:peps_czx}

In the CZX model of \cite{ChenLiuWen2011CZX} every site of the square lattice
carries four qubits, one at each corner, and the ground state on a closed
surface is the product over plaquettes of the states $\ket{0000}+\ket{1111}$
on the four qubits around each plaquette. Blocking the four qubits of a site,
\cite[Appendix~A, ``The CZX model'']{Cirac2021Matrix} writes this state as a
PEPS of bond dimension four.

\begin{definition}[CZX tensor]%
    \label{def:peps_czx_tensor}
    \lean{TNLean.PEPS.czxBond, TNLean.PEPS.czxQubits, TNLean.PEPS.czxTopLeft,
        TNLean.PEPS.czxTopRight, TNLean.PEPS.czxBottomRight,
        TNLean.PEPS.czxBottomLeft, TNLean.PEPS.czxSiteTensor,
        TNLean.PEPS.czxBondSwap, TNLean.PEPS.czxBondSwap_czxBond,
        TNLean.PEPS.czxPEPS}
    \leanok
    \uses{def:peps_torus_site_tensor}
    Label the four qubits of a site $\ket{ijkl}$, with $i$ at the top left,
    $j$ at the top right, $k$ at the bottom right and $l$ at the bottom left,
    and label the four basis vectors of a bond by pairs of qubits. The CZX
    tensor is
    \begin{align}
        \sum_{i,j,k,l=0}^{1}\ket{ijkl}\bigl((i,j),(j,k),(k,l),(l,i)\bigr|,
        \notag
    \end{align}
    so each leg carries the pair of corner qubits on its side, listed
    clockwise around the site
    \cite[Appendix~A, ``The CZX model'']{Cirac2021Matrix}. The CZX PEPS
    places this tensor at every site with its down and left labels read in
    reverse, so that each bond identifies the pair $(a,b)$ at one end with
    the pair $(b,a)$ at the other.
\end{definition}

The reversal corrects the printed formula. The review contracts every bond
with equal labels at its two ends
\cite[Section~``MPS and PEPS'']{Cirac2021Matrix}. Contracted in this way,
the printed tensor sets the top-right qubit of a site equal to the
bottom-left qubits of its right and upper neighbours, and the resulting state
is a product of GHZ states on diagonal loops that wind around the torus, not
the plaquette state. Two neighbouring sites list the qubits of their common
side in opposite orders, and the plaquette state is obtained by identifying
the pair $(a,b)$ at one end of a bond with $(b,a)$ at the other
\cite{gap:rmp_peps_czx_bond_orientation}.

\begin{theorem}[The CZX PEPS is a product of plaquette GHZ states]%
    \label{thm:peps_czx_plaquette_product}
    \lean{TNLean.PEPS.czx_bond_iff_plaquette, TNLean.PEPS.stateCoeff_czxPEPS}
    \leanok
    \uses{def:peps_czx_tensor, def:peps_stateCoeff}
    On a torus with $w,h\ge 3$ the CZX PEPS has coefficient
    \begin{align}
        \Coeff(\sigma)=\prod_{v}\delta\bigl[j_v=i_{v+(1,0)}=k_{v+(0,1)}
            =l_{v+(1,1)}\bigr],
        \notag
    \end{align}
    where $(i_v,j_v,k_v,l_v)$ are the four qubits of $\sigma(v)$ and the
    factor at $v$ is the coefficient of $\ket{0000}+\ket{1111}$ on the four
    qubits around the plaquette at the upper right corner of $v$
    \cite[Appendix~A, ``The CZX model'']{Cirac2021Matrix},
    \cite{ChenLiuWen2011CZX}.
\end{theorem}

\begin{proof}\leanok
    \uses{lem:peps_torus_site_tensor_contraction}
    By Lemma~\ref{lem:peps_torus_site_tensor_contraction} the top and right
    legs force the vertical bond above $v$ to carry $(i_v,j_v)$ and the
    horizontal bond to its right to carry $(j_v,k_v)$. The down and left legs
    then require $(i_{v-(0,1)},j_{v-(0,1)})=(l_v,k_v)$ and
    $(j_{v-(1,0)},k_{v-(1,0)})=(i_v,l_v)$ at every site, and these
    conditions, shifted by one site, are equivalent to the plaquette
    conditions.
\end{proof}

The on-site $\mathbb{Z}_2$ symmetry of the CZX model is generated by
$U_{CZX}=U_XU_{CZ}$, where $U_X=X_1X_2X_3X_4$ flips the four qubits of a site
and $U_{CZ}=CZ_{12}CZ_{23}CZ_{34}CZ_{41}$ applies controlled-$Z$ gates to its
four pairs of neighbouring qubits \cite{ChenLiuWen2011CZX}. On the state
$\ket{ijkl}$ it acts by
$U_{CZX}\ket{ijkl}=(-1)^{ij+jk+kl+li}\ket{\bar\imath\bar\jmath\bar k\bar l}$.

\begin{definition}[CZX on-site symmetry]%
    \label{def:peps_czx_on_site}
    \lean{TNLean.PEPS.onSiteOperator, TNLean.PEPS.onSiteOperator_apply,
        TNLean.PEPS.onSiteOperator_monomial, TNLean.PEPS.czxSiteFlip,
        TNLean.PEPS.czxTopLeft_czxSiteFlip, TNLean.PEPS.czxTopRight_czxSiteFlip,
        TNLean.PEPS.czxBottomRight_czxSiteFlip,
        TNLean.PEPS.czxBottomLeft_czxSiteFlip, TNLean.PEPS.czxSiteFlip_czxQubits,
        TNLean.PEPS.czxSiteFlip_czxSiteFlip, TNLean.PEPS.czxSiteFlip_symm,
        TNLean.PEPS.czxCZExponent, TNLean.PEPS.czxOnSite,
        TNLean.PEPS.czxOnSite_eq_monomial}
    \leanok
    \uses{def:peps_czx_tensor}
    The operator $U_{CZX}=U_XU_{CZ}$ on the sixteen states of a site, and
    its product $U_{CZX}^{\otimes V}$ over the sites $V$ of the lattice, whose
    entry at $(\sigma',\sigma)$ is $\prod_v(U_{CZX})_{\sigma'_v\sigma_v}$
    \cite{ChenLiuWen2011CZX}.
\end{definition}

\begin{theorem}[The CZX PEPS is invariant under the on-site symmetry]%
    \label{thm:peps_czx_invariant}
    \lean{TNLean.PEPS.prod_neg_one_pow_czxCZExponent,
        TNLean.PEPS.onSiteOperator_czxOnSite_mulVec_stateCoeff_czxPEPS}
    \leanok
    \uses{def:peps_czx_on_site, thm:peps_czx_plaquette_product}
    On a torus with $w,h\ge 3$ the CZX PEPS is invariant under
    $U_{CZX}^{\otimes V}$
    \cite[Appendix~A, ``The CZX model'']{Cirac2021Matrix},
    \cite{ChenLiuWen2011CZX}.
\end{theorem}

\begin{proof}\leanok
    \uses{thm:peps_czx_plaquette_product}
    Flipping every qubit preserves the plaquette constraints of
    Theorem~\ref{thm:peps_czx_plaquette_product}. On a configuration
    satisfying them, let $P_v$ be the value of the plaquette at the upper
    right corner of $v$. The sign of $U_{CZX}$ at $v$ is
    $(-1)^{e_v}$ with
    $e_v=P_{v-(1,0)}P_v+P_vP_{v-(0,1)}+P_{v-(0,1)}P_{v-(1,1)}
        +P_{v-(1,1)}P_{v-(1,0)}$. Shifting the summation variable, the sums
    over $v$ of the third and fourth terms equal those of the first and
    second, so $\sum_ve_v$ is even: each controlled-$Z$ gate between two
    neighbouring plaquettes appears twice.
\end{proof}

\begin{definition}[CZX leg operator]%
    \label{def:peps_czx_leg_operator}
    \lean{TNLean.PEPS.czxLegFlip, TNLean.PEPS.czxLegFlip_czxBond,
        TNLean.PEPS.czxLegPhase, TNLean.PEPS.czxLegPhase_czxBond,
        TNLean.PEPS.czxLegOperator}
    \leanok
    \uses{def:peps_czx_tensor}
    A virtual leg carries two qubits $(a,b)$. The operator
    $V=(X\otimes X)\,CZ$ acts on a leg by
    $V\ket{ab}=(-1)^{ab}\ket{\bar a\bar b}$.
\end{definition}

\begin{theorem}[Pulling the symmetry through a site]%
    \label{thm:peps_czx_pull_through}
    \lean{TNLean.PEPS.czxOnSite_mul_czxSiteTensor}
    \leanok
    \uses{def:peps_czx_on_site, def:peps_czx_leg_operator}
    Applying $U_{CZX}$ to the physical index of the CZX tensor equals applying
    $V$ to each of its four virtual legs
    \cite[Section~III.B, ``Non-injective PEPS: SPT phases'']{Cirac2021Matrix}:
    writing $A^s_{trbl}$ for the entry of the CZX tensor with physical index $s$
    and top, right, bottom and left legs $t,r,b,l$,
    \begin{align}
        \sum_s(U_{CZX})_{s's}\,A^s_{trbl}
         & =\sum_{t',r',b',l'}A^{s'}_{t'r'b'l'}\,V_{t't}V_{r'r}V_{b'b}V_{l'l}.
        \label{eq:peps_czx_pull_through}
    \end{align}
\end{theorem}

\begin{proof}\leanok
    Both sides are supported on the flipped legs of the flipped qubits, and
    the sign $(-1)^{ij+jk+kl+li}$ of $U_{CZX}$ is the product of the signs
    of $V$ on the four legs $(i,j)$, $(j,k)$, $(k,l)$, $(l,i)$.
\end{proof}

On the boundary of a region, neighbouring boundary legs carry the same qubit
of their common plaquette, one effective spin per boundary plaquette
\cite{ChenLiuWen2011CZX}.

\begin{definition}[Boundary legs of the CZX tensor network]%
    \label{def:peps_czx_boundary_embedding}
    \lean{TNLean.PEPS.czxBoundaryLegs, TNLean.PEPS.czxBoundaryEmbedding}
    \leanok
    \uses{def:peps_czx_tensor}
    For $N>0$ and effective spins $c\in\{0,1\}^N$, the closed chain of $N$
    boundary legs carrying $c$ has the pair $(c_m,c_{m+1})$ on leg $m$, indices
    modulo $N$. The map $\iota$ from $N$ effective spins into the legs is
    \begin{align}
        \iota\ket{c} & =\ket{(c_0,c_1),(c_1,c_2),\dots,(c_{N-1},c_0)}.
        \label{eq:peps_czx_boundary_embedding}
    \end{align}
\end{definition}

\begin{theorem}[The boundary action is the CZX matrix product unitary]%
    \label{thm:peps_czx_boundary_mpu}
    \lean{TNLean.PEPS.czxBoundaryLegs_injective,
        TNLean.PEPS.czxBoundaryEmbedding_conjTranspose_mul_self,
        TNLean.PEPS.onSiteOperator_czxLegOperator_mul_czxBoundaryEmbedding,
        TNLean.PEPS.czxBoundaryEmbedding_conjTranspose_mul_onSiteOperator_mul,
        TNLean.PEPS.onSiteOperator_czxLegOperator_mul_czxBoundaryEmbedding_review}
    \leanok
    \uses{def:peps_czx_leg_operator, def:peps_czx_boundary_embedding,
        thm:asymex_czx_kernel, def:asymex_czx_review_tensor}
    For every $N>0$, the chain of boundary legs determines $c$, so $\iota$ is an
    isometry, $\iota^\dagger\iota=\Id$. It satisfies $V^{\otimes N}\iota=\iota\,U_N$ and
    $\iota^\dagger V^{\otimes N}\iota=U_N$, where
    $U_N\ket{c}=(-1)^{\sum_mc_mc_{m+1}}\ket{\bar c}$ is the CZX matrix product
    unitary of Theorem~\ref{thm:asymex_czx_kernel}; in terms of the
    review's operator $O(A)$ with the printed matrices,
    $V^{\otimes N}\iota=(-1)^N\iota\,O(A)$
    \cite[Appendix~A, ``The CZX MPU'']{Cirac2021Matrix},
    \cite{ChenLiuWen2011CZX}. For actual positive proper bounded rectangles, the physical region image
    and this cyclic boundary action are identified in
    Section~\ref{sec:peps_czx_rectangle_boundary}, together with the actual
    reduced-density support and rank, preserved by trace normalization.
    Nonzero spectral values, entropy, isometric normalization of the physical
    boundary map, and nonrectangular regions remain outside that result
    \cite{gap:rmp_peps_czx_boundary_chain}.
\end{theorem}

\begin{proof}\leanok
    \uses{thm:asymex_czx_kernel, thm:asymex_czx_review_kernel}
    $V^{\otimes N}$ maps the legs of $c$ to the legs of $\bar c$ with the sign
    $\prod_m(-1)^{c_mc_{m+1}}$, which is the kernel of
    Theorem~\ref{thm:asymex_czx_kernel}; the relation to $O(A)$ is
    Theorem~\ref{thm:asymex_czx_review_kernel}.
\end{proof}

The review states that the CZX model lies in the non-trivial
symmetry-protected sector of its on-site $\mathbb{Z}_2$ symmetry
\cite[Appendix~A, ``The CZX model'']{Cirac2021Matrix}. This classification is
not formalized for the two-dimensional state. Its boundary symmetry is a
matrix product unitary; the three-cocycle attached to a group of matrix
product unitaries is constructed in
Theorems~\ref{thm:mpug_anomaly_three_cocycle}
and~\ref{thm:mpug_anomaly_three_cocycle_gauge}, and for the CZX boundary
representation its class is computed and shown to be non-trivial in
Theorem~\ref{thm:asymex_czx_decorated_anomaly_class}.

\section{The physical boundary symmetry of a CZX rectangle}%
\label{sec:peps_czx_rectangle_boundary}

Let $R=[x_0,x_0+w)\times[y_0,y_0+h)$ be a coordinate rectangle in a
$W\times H$ torus. Assume $W,H\geq3$, $w,h>0$, $w<W$, $h<H$, and
$x_0+w\leq W$, $y_0+h\leq H$. Write $P=2w+2h$.
The following construction uses the actual open-region map $A_R$, summing
only internal bonds and leaving crossing bonds free. It identifies the
plaquette spins in the boundary discussion of \cite{ChenLiuWen2011CZX},
source lines 330--345, with the image of this physical region map.

\begin{definition}[Clockwise native rectangle boundary]%
    \label{def:czx_rectangle_perimeter}
    \lean{TNLean.PEPS.rectanglePerimeterCorner,
        TNLean.PEPS.rectanglePerimeterEnd,
        TNLean.PEPS.rectanglePerimeterReversed,
        TNLean.PEPS.torusRectanglePerimeterCode,
        TNLean.PEPS.torusRectanglePerimeterEdge,
        TNLean.PEPS.torusRectanglePerimeterSite,
        TNLean.PEPS.torusRectanglePerimeterLeg,
        TNLean.PEPS.torusRectanglePerimeterCorner}
    \leanok
    \uses{def:peps_torus_rectangular_grid}
    Number the crossing bonds clockwise, first across the top, then down
    the right, across the bottom, and up the left. The corresponding first
    plaquette corners are
    $(i,h)$, $(w,h-j)$, $(w-k,0)$, and $(0,\ell)$ on the four sides.
    Reverse the native pair order on the bottom and left.
\end{definition}

\begin{theorem}[Native perimeter and distinct plaquette corners]%
    \label{thm:czx_rectangle_perimeter}
    \lean{TNLean.PEPS.even_rectanglePerimeter,
        TNLean.PEPS.rectanglePerimeterEnd_eq_corner_add_one,
        TNLean.PEPS.torusRectanglePerimeterEdge_boundary,
        TNLean.PEPS.torusRectanglePerimeterEdge_injective,
        TNLean.PEPS.torusRectangleBoundaryEquiv,
        TNLean.PEPS.torusRectangleBoundaryEquiv_apply_val,
        TNLean.PEPS.torusRectanglePerimeterSite_mem,
        TNLean.PEPS.torusRectanglePerimeterLeg_val,
        TNLean.PEPS.rectanglePerimeterCorner_le,
        TNLean.PEPS.rectanglePerimeterCorner_injective,
        TNLean.PEPS.torusRectanglePerimeterCorner_injective}
    \leanok
    \uses{def:czx_rectangle_perimeter}
    This numbering is a bijection from $\{0,\ldots,P-1\}$ to the actual
    crossing bonds. Its cyclic successor is ordinary addition in
    $\mathbb Z/P\mathbb Z$. Every bond has the specified endpoint in $R$,
    the perimeter corners are distinct in the torus, and $P$ is even.
    The assertions include rectangles meeting a coordinate seam.
\end{theorem}
\begin{proof}\leanok
    \uses{lem:peps_torus_rectangle_boundary_card}
    The four native edge families are disjoint and injective. The existing
    crossing-bond cardinality theorem gives surjectivity. Proper side
    lengths make the closed corner grid injective modulo each torus period.
\end{proof}

\begin{theorem}[Exact coefficients and disjoint boundary columns]%
    \label{thm:czx_open_region_coefficients}
    \lean{TNLean.PEPS.czx_regionIncidentConfig_unique,
        TNLean.PEPS.openRegionWeight_czxPEPS,
        TNLean.PEPS.openRegionWeight_czxPEPS_disjoint,
        TNLean.PEPS.openRegionCoefficientMatrix_czxPEPS_gram_offDiagonal}
    \leanok
    \uses{def:peps_open_region_contraction,def:peps_czx_tensor}
    For any finite region, a physical configuration determines at most one
    nonzero incident-bond assignment. Its open-region coefficient is one
    precisely when this assignment exists and has the specified boundary;
    otherwise it is zero. Different boundary columns have disjoint physical
    supports, so the region Gram matrix is diagonal in the native boundary
    basis. No labels on edges wholly outside the region are summed.
\end{theorem}
\begin{proof}\leanok
    The printed CZX tensor fixes all four incident labels from the four
    physical qubits. Every incident edge touches a region vertex, so this
    determines the entire assignment. The contraction sum has at most one
    nonzero summand, whose value is one.
\end{proof}

\begin{definition}[Actual crossing-bond symmetry]%
    \label{def:czx_region_boundary_symmetry}
    \lean{TNLean.PEPS.czxRegionBoundaryFlip,
        TNLean.PEPS.czxRegionBoundaryPhase}
    \leanok
    \uses{def:peps_czx_leg_operator,def:peps_open_region_contraction}
    For a native boundary configuration $\mu$, let $f_R\mu$ flip both bits
    on each crossing bond and let
    $\phi_R(\mu)=\prod_{e\in\partial R}(-1)^{a_eb_e}$, where
    $\mu_e=(a_e,b_e)$.
\end{definition}

\begin{theorem}[Physical region pull-through]%
    \label{thm:czx_region_physical_pullthrough}
    \lean{TNLean.PEPS.czxLegPhase_swap,
        TNLean.PEPS.czxLegFlip_swap,
        TNLean.PEPS.regionPhysicalMap_czxOnSite_openRegionWeight,
        TNLean.PEPS.regionPhysicalProductMatrix_czxOnSite_intertwines}
    \leanok
    \uses{def:czx_region_boundary_symmetry,def:peps_czx_on_site}
    Let $U_R$ apply the CZX on-site symmetry at each site of $R$. Then
    \[
        U_RA_R|\mu\rangle=\phi_R(\mu)A_R|f_R\mu\rangle.
    \]
    This identity holds for the actual open contraction of every finite
    region. Reversal of a native bond pair preserves its phase and commutes
    with its simultaneous bit flip.
\end{theorem}
\begin{proof}\leanok
    \uses{thm:peps_czx_pull_through,thm:peps_region_physical_maps}
    Apply the one-site pull-through identity inside the finite contraction.
    Reindex vertex--edge incidences by edge--endpoint incidences. Each
    internal bond contributes the square of its controlled phase, hence
    one; each crossing bond contributes its phase once. Reindex the
    incident-bond sum by simultaneous bit flips.
\end{proof}

\begin{definition}[Oriented rectangle boundary labels]%
    \label{def:czx_rectangle_boundary_labels}
    \lean{TNLean.PEPS.czxRectangleBoundaryLabels,
        TNLean.PEPS.czxRectangleCoordinatesEquiv,
        TNLean.PEPS.czxRectangleCoordinatesEquiv_apply,
        TNLean.PEPS.czxRectangleEffectiveBoundaryConfig,
        TNLean.PEPS.czxRectangleBoundaryMatrix}
    \leanok
    \uses{def:czx_rectangle_perimeter,def:peps_czx_boundary_embedding,
        def:peps_open_region_contraction}
    Let $E_R\mu$ be the clockwise label sequence, with bottom and left
    pairs reversed. This is a bijective change of coordinates. For
    $c\in\{0,1\}^P$, set $\iota(c)_i=(c_i,c_{i+1})$, define
    $\mu(c)=E_R^{-1}\iota(c)$, and let $F_R|c\rangle=A_R|\mu(c)\rangle$.
\end{definition}

\begin{theorem}[The contraction imposes the cyclic plaquette constraint]%
    \label{thm:czx_rectangle_boundary_support}
    \lean{TNLean.PEPS.czxRectangleBoundaryLabels_mem_range_of_openRegionWeight_ne_zero}
    \leanok
    \uses{def:czx_rectangle_boundary_labels}
    If a coefficient of $A_R|\mu\rangle$ is nonzero, then
    $E_R\mu=\iota(c)$ for some $c\in\{0,1\}^P$.
\end{theorem}
\begin{proof}\leanok
    \uses{thm:czx_rectangle_perimeter,thm:czx_open_region_coefficients}
    Extract the nonzero incident-bond assignment. On each straight side,
    the shared internal bond equates the adjacent physical qubits. At each
    corner, the equality is within one site tensor. These eight cases give
    the cyclic adjacent-pair constraint, including the final wraparound.
\end{proof}

\begin{definition}[Plaquette-bit witnesses]%
    \label{def:czx_plaquette_witness}
    \lean{TNLean.PEPS.czxPlaquettePhysical,TNLean.PEPS.czxPlaquetteBondConfig}
    \leanok
    \uses{def:peps_czx_tensor}
    Given bits $q$ at torus plaquette corners, assign site $(x,y)$ the qubits
    $(q_{x,y+1},q_{x+1,y+1},q_{x+1,y},q_{x,y})$ and give each native bond
    the corresponding adjacent pair.
\end{definition}

\begin{theorem}[Exact witnesses and a faithful effective boundary]%
    \label{thm:czx_rectangle_boundary_faithful}
    \lean{TNLean.PEPS.czxPlaquetteBondConfig_right,
        TNLean.PEPS.czxPlaquetteBondConfig_up,
        TNLean.PEPS.component_czxPlaquettePhysical,
        TNLean.PEPS.openRegionWeight_czxPlaquettePhysical,
        TNLean.PEPS.czxRectangleBoundaryLabels_plaquette,
        TNLean.PEPS.czxRectangleEffectiveBoundaryConfig_injective,
        TNLean.PEPS.exists_czxRectangleBoundaryMatrix_selector,
        TNLean.PEPS.czxRectangleBoundaryMatrix_mulVec_injective}
    \leanok
    \uses{def:czx_rectangle_boundary_labels,def:czx_plaquette_witness}
    Each local tensor and each actual region witness coefficient equals
    one. For every $c$ there is a physical configuration $\sigma_c$ with
    $F_R(\sigma_c,c')=\delta_{c,c'}$. In particular, $F_R$ is injective.
\end{theorem}
\begin{proof}\leanok
    \uses{thm:czx_rectangle_perimeter,thm:czx_open_region_coefficients}
    Extend arbitrary boundary bits from the distinct perimeter corners to
    torus plaquette bits, assigning zero elsewhere. The resulting physical
    configuration has coefficient one in the selected column and zero in
    all other columns by disjointness.
\end{proof}

\begin{theorem}[Exact physical image and its dimension]%
    \label{thm:czx_rectangle_physical_image}
    \lean{TNLean.PEPS.range_openRegionMap_czxRectangle,
        TNLean.PEPS.finrank_range_openRegionMap_czxRectangle}
    \leanok
    \uses{def:czx_rectangle_boundary_labels}
    The image of the actual region map is exactly the span of the columns
    of $F_R$, and its dimension is $2^P$.
\end{theorem}
\begin{proof}\leanok
    \uses{thm:czx_rectangle_boundary_support,thm:czx_rectangle_boundary_faithful}
    Every nonzero native boundary column is an effective column; every
    effective column is a native one. The selectors prove independence,
    and there are $2^P$ effective spin configurations.
\end{proof}

\begin{theorem}[Physical action on the effective plaquette spins]%
    \label{thm:czx_rectangle_effective_symmetry}
    \lean{TNLean.PEPS.regionPhysicalMap_czxRectangleBoundaryMatrix_column}
    \leanok
    \uses{def:czx_rectangle_boundary_labels}
    On the faithful physical boundary coordinates,
    \[
        U_RF_R|c\rangle
        =(-1)^{\sum_i c_ic_{i+1}}F_R|\bar c\rangle.
    \]
\end{theorem}
\begin{proof}\leanok
    \uses{thm:czx_region_physical_pullthrough,thm:czx_rectangle_boundary_faithful}
    Reorder crossing bonds clockwise. Pair reversal preserves the controlled
    phase and commutes with bit flip, giving one controlled phase for each
    adjacent pair of effective spins.
\end{proof}

\begin{theorem}[The printed matrix product operator on the physical rectangle image]%
    \label{thm:czx_rectangle_review_mpo}
    \lean{TNLean.PEPS.regionPhysicalProductMatrix_czxRectangleBoundaryMatrix_review}
    \uses{def:czx_rectangle_boundary_labels,def:asymex_czx_review_tensor}
    The review's printed CZX matrix product operator $O(A)$ satisfies
    $U_RF_R=F_RO(A)$ on the actual rectangle image.
\end{theorem}
\begin{proof}
    \uses{thm:czx_rectangle_effective_symmetry,thm:czx_rectangle_perimeter,
        thm:asymex_czx_kernel,thm:asymex_czx_review_kernel}
    The printed matrices give $D_PX^{\otimes P}=(-1)^PX^{\otimes P}D_P$.
    The rectangle perimeter $P$ is even, so this sign is one. The normalized
    bra-ket display is not substituted for those printed matrices.
\end{proof}

The identification above is a linear identification of the actual
open-region image. No formula for the normalized reduced density matrix,
its flatness, or its entropy is asserted here. Nonrectangular regions
remain outside the perimeter construction.


\section{The AKLT state}%
\label{sec:peps_aklt}

The two-dimensional AKLT state places singlets on the edges of the lattice and
projects the four spin-$\tfrac12$ at each site onto their symmetric subspace
\cite[Appendix~A, ``The AKLT model'']{Cirac2021Matrix}. Label the basis of
four spin-$\tfrac12$ by $s=(s_0,s_1,s_2,s_3)\in\{0,1\}^4$, write $|s|$ for the
number of entries equal to $1$, and let $P_{\pi}$ be the permutation of the four
factors by $\pi\in S_4$.

\begin{definition}[Symmetric projector on four spin-$\tfrac12$]%
    \label{def:peps_aklt_sym_projector}
    \lean{TNLean.PEPS.legKronecker, TNLean.PEPS.legPermMatrix,
        TNLean.PEPS.symProjector}
    \leanok
    The projector onto the symmetric subspace of $(\C^2)^{\otimes 4}$ is
    $\Pi_{\mathrm{sym}}=\frac{1}{24}\sum_{\pi\in S_4}P_{\pi}$.
\end{definition}

\begin{lemma}[The symmetric subspace has dimension five]%
    \label{lem:peps_aklt_sym_projector}
    \lean{TNLean.PEPS.legKronecker_mul, TNLean.PEPS.legKronecker_one,
        TNLean.PEPS.legKronecker_transpose,
        TNLean.PEPS.legPermMatrix_mul, TNLean.PEPS.legPermMatrix_conjTranspose,
        TNLean.PEPS.legPermMatrix_mul_legKronecker,
        TNLean.PEPS.symProjector_mul_self, TNLean.PEPS.symProjector_conjTranspose,
        TNLean.PEPS.symProjector_mul_legKronecker, TNLean.PEPS.legWeight,
        TNLean.PEPS.card_perm_comp_eq, TNLean.PEPS.symProjector_apply,
        TNLean.PEPS.choose_legWeight_dvd, TNLean.PEPS.sum_div_choose_legWeight,
        TNLean.PEPS.trace_symProjector, TNLean.PEPS.finrank_range_symProjector}
    \leanok
    \uses{def:peps_aklt_sym_projector}
    $\Pi_{\mathrm{sym}}$ is a Hermitian idempotent commuting with $U^{\otimes 4}$
    for every $2\times 2$ matrix $U$. Its entries are
    $(\Pi_{\mathrm{sym}})_{s,t}=\binom{4}{|s|}^{-1}$ if $|s|=|t|$ and $0$
    otherwise, and its range, the spin-$2$ subspace, has dimension five.
\end{lemma}

\begin{proof}\leanok
    $P_{\pi}P_{\rho}=P_{\rho\pi}$ and $P_{\pi}^{\dagger}=P_{\pi^{-1}}$, so the average
    over $S_4$ is a Hermitian idempotent, and
    $P_{\pi} U^{\otimes 4}=U^{\otimes 4}P_{\pi}$ by reordering the tensor
    factors. The number of $\pi$ with
    $s=t\circ\pi$ is $24/\binom{4}{|t|}$ when $|s|=|t|$ and $0$ otherwise, by
    a finite check. The dimension of the range of an idempotent is its trace,
    $\sum_{s}\binom{4}{|s|}^{-1}=\sum_{k=0}^{4}\binom4k\binom4k^{-1}=5$.
\end{proof}

\begin{definition}[AKLT tensor]%
    \label{def:peps_aklt_tensor}
    \lean{TNLean.PEPS.akltLegs,
        TNLean.PEPS.akltSiteTensorFun, TNLean.PEPS.akltSiteTensor,
        TNLean.PEPS.akltPEPS}
    \leanok
    \uses{def:peps_aklt_sym_projector, def:peps_torus_site_tensor,
        def:aklt_tensor_review}
    With the singlet matrix
    $Y=\left(\begin{smallmatrix}0&-1\\1&0\end{smallmatrix}\right)$, the AKLT
    tensor on the square lattice is
    $A=\Pi_{\mathrm{sym}}(\mathbb{1}\otimes\mathbb{1}\otimes Y\otimes Y)$,
    with physical index $s\in\{0,1\}^4$ (physical dimension $16$) and virtual
    indices ordered top, right, down, left
    \cite[Appendix~A, ``The AKLT model'']{Cirac2021Matrix}.
\end{definition}

\begin{lemma}[Entries and physical span of the AKLT tensor]%
    \label{lem:peps_aklt_entries}
    \lean{Matrix.singletY_apply, Matrix.singletY_mul_transpose,
        TNLean.PEPS.akltLegs_apply,
        TNLean.PEPS.legWeight_vec, TNLean.PEPS.akltSiteTensorFun_apply,
        TNLean.PEPS.akltLegs_mul_transpose, TNLean.PEPS.span_akltSiteTensorFun}
    \leanok
    \uses{def:peps_aklt_tensor, lem:peps_aklt_sym_projector}
    The entries of the AKLT tensor are
    \begin{align}
        A^{s}_{\alpha\beta\gamma\delta}=
        \frac{(-1)^{\gamma+\delta}}{\binom{4}{k}}\,\delta_{|s|,k},\qquad
        k=\alpha+\beta+(1-\gamma)+(1-\delta),
        \notag
    \end{align}
    and the span of the vectors $s\mapsto A^s_{\alpha\beta\gamma\delta}$ over
    all virtual indices is the range of $\Pi_{\mathrm{sym}}$.
\end{lemma}

\begin{proof}\leanok
    $Y\ket{0}=\ket{1}$ and $Y\ket{1}=-\ket{0}$, so the column
    $(\alpha,\beta,\gamma,\delta)$ of $\mathbb{1}\otimes\mathbb{1}\otimes
        Y\otimes Y$ is $(-1)^{\gamma+\delta}$ times the basis vector
    $(\alpha,\beta,1-\gamma,1-\delta)$. The matrix
    $L=\mathbb{1}\otimes\mathbb{1}\otimes Y\otimes Y$ satisfies $LL^{T}=\mathbb{1}$,
    so the columns of $\Pi_{\mathrm{sym}}L$ span the range of
    $\Pi_{\mathrm{sym}}$.
\end{proof}

\begin{theorem}[$SU(2)$ covariance of the AKLT tensor]%
    \label{thm:peps_aklt_su2}
    \lean{Matrix.mul_singletY, TNLean.PEPS.akltSiteTensor_covariant,
        Matrix.adjugate_eq_star_of_mem_specialUnitaryGroup,
        TNLean.PEPS.akltSiteTensor_su2_covariant}
    \leanok
    \uses{def:peps_aklt_tensor, lem:peps_aklt_sym_projector}
    Regard $A$ as the map from the virtual space $(\C^2)^{\otimes 4}$ to the
    physical space. For every $U\in SU(2)$,
    \begin{align}
        U^{\otimes 4}A=A\,(U\otimes U\otimes\bar U\otimes\bar U).
        \notag
    \end{align}
    For an arbitrary $2\times 2$ matrix $U$ the same identity holds with $\bar U$ replaced by
    $(\operatorname{adj}U)^{T}$.
\end{theorem}

\begin{proof}\leanok
    For a $2\times2$ matrix, $UY=Y(\operatorname{adj}U)^{T}$ by direct
    computation. Since $\Pi_{\mathrm{sym}}$ commutes with $U^{\otimes 4}$,
    $U^{\otimes4}\Pi_{\mathrm{sym}}L=\Pi_{\mathrm{sym}}(U\otimes U\otimes
        UY\otimes UY)=\Pi_{\mathrm{sym}}L\,(U\otimes U\otimes V\otimes V)$ with
    $V=(\operatorname{adj}U)^{T}$. For $U\in SU(2)$, $U^{\dagger}U=\mathbb{1}$
    and $U\operatorname{adj}U=(\det U)\mathbb{1}=\mathbb{1}$, so
    $\operatorname{adj}U=U^{\dagger}U\operatorname{adj}U=U^{\dagger}$ and
    $V=\bar U$.
\end{proof}

\begin{theorem}[The AKLT PEPS as projected singlets]%
    \label{thm:peps_aklt_state}
    \lean{TNLean.PEPS.akltLegs_vec, TNLean.PEPS.stateCoeff_akltPEPS}
    \leanok
    \uses{def:peps_aklt_tensor, lem:peps_torus_site_tensor_contraction}
    On a torus with $w,h\ge 3$ the AKLT PEPS has coefficient
    \begin{align}
        \Coeff(\sigma)=\sum_{\tau}\prod_{v}(\Pi_{\mathrm{sym}})_{\sigma_v,\tau_v}
        \prod_{v}Y_{\tau_{v+(1,0),3},\,\tau_{v,1}}\,Y_{\tau_{v+(0,1),2},\,\tau_{v,0}},
        \notag
    \end{align}
    the sum running over the spin-$\tfrac12$ labels $\tau_v\in\{0,1\}^4$ of the
    four legs of every site. The factor $Y_{b,a}$ is the coefficient of
    $\ket{a}\ket{b}$ in the singlet $\ket{01}-\ket{10}$, so the state is
    $\bigotimes_v\Pi_{\mathrm{sym}}$ applied to a singlet on every edge.
\end{theorem}

\begin{proof}\leanok
    \uses{lem:peps_torus_site_tensor_contraction}
    Expand each $A=\Pi_{\mathrm{sym}}L$ as a sum over the intermediate label
    $\tau_v$. The identity factors of $L$ force the bond index on the top edge
    of $v$ to be $\tau_{v,0}$ and on the right edge to be $\tau_{v,1}$, and the
    $Y$ factors on the down and left legs then give
    $Y_{\tau_{v,2},\tau_{v-(0,1),0}}Y_{\tau_{v,3},\tau_{v-(1,0),1}}$; shifting
    the product by one site gives the formula.
\end{proof}

The tensor $A$ takes values in $(\C^2)^{\otimes 4}$ with physical support on
the five-dimensional symmetric subspace, on which the physical action
$U^{\otimes 4}$ acts as the irreducible spin-$2$ representation
\cite[Appendix~A, ``The AKLT model'']{Cirac2021Matrix}.

\section{The RVB state}%
\label{sec:peps_rvb}

The nearest-neighbour resonating valence bond state is the superposition of
all coverings of the lattice by nearest-neighbour singlets. Its PEPS of bond
dimension three combines, at each site, the projector
\begin{align}
    P=\ket{0}\bigl[(0222|+(2022|+(2202|+(2220|\bigr]
    +\ket{1}\bigl[(1222|+(2122|+(2212|+(2221|\bigr]
    \notag
\end{align}
with singlet matrices on the edges; on the square lattice with a
translation-invariant orientation of the singlets the tensor is
$A=P(\mathbb{1}\otimes\mathbb{1}\otimes Y\otimes Y)$
\cite[Appendix~A, ``The RVB state'']{Cirac2021Matrix}. On the bond of
dimension three the matrix $Y$ acts on the spin states $0,1$ and is extended by
$1$ on the vacuum state $2$, so that the bond state is
$\ket{01}-\ket{10}+\ket{22}$. The extension by $0$ would make every entry of
$A$ vanish, since a leg carrying $Y$ could then never be in the vacuum state
\cite{gap:rmp_peps_rvb_bond_vacuum}.

\begin{definition}[RVB tensor]%
    \label{def:peps_rvb_tensor}
    \lean{TNLean.PEPS.rvbSingletY, TNLean.PEPS.rvbProjector,
        TNLean.PEPS.rvbSiteTensorInt, TNLean.PEPS.rvbSiteTensor,
        TNLean.PEPS.rvbPEPS}
    \leanok
    \uses{def:peps_torus_site_tensor}
    The RVB tensor is $A=P(\mathbb{1}\otimes\mathbb{1}\otimes Y\otimes Y)$ with
    $Y=\left(\begin{smallmatrix}0&-1&0\\1&0&0\\0&0&1\end{smallmatrix}\right)$,
    virtual indices ordered top, right, down, left.
\end{definition}

\begin{lemma}[Local dimer constraint]%
    \label{lem:peps_rvb_local}
    \lean{TNLean.PEPS.rvbFlip, TNLean.PEPS.rvbSiteTensorInt_eq,
        TNLean.PEPS.rvbSiteTensor_apply, TNLean.PEPS.rvbSiteTensor_ne_zero_iff,
        TNLean.PEPS.rvbBondLabel, TNLean.PEPS.rvbBondLabel_ne_two_iff,
        TNLean.PEPS.rvbSiteTensorInt_ne_zero,
        TNLean.PEPS.rvbSiteTensorInt_bondLabel}
    \leanok
    \uses{def:peps_rvb_tensor}
    $A^{s}_{urdl}\neq 0$ if and only if exactly one of $u,r,d,l$ differs from
    $2$, that index being $s$ if it is $u$ or $r$ and $1-s$ if it is $d$ or
    $l$. The nonzero entries equal $-1$ when the down or left index is $1$ and
    $+1$ otherwise.
\end{lemma}

\begin{proof}\leanok
    $Y$ maps the bond index $d$ to $1-d$ on the spin states, with sign $-1$
    at $d=1$, and fixes the vacuum. The statement is a finite check over the
    $162$ entries.
\end{proof}

\begin{theorem}[The RVB PEPS is the superposition of singlet coverings]%
    \label{thm:peps_rvb_state}
    \lean{TNLean.PEPS.IsTorusDimerCovering, TNLean.PEPS.rvbCoveringBonds,
        TNLean.PEPS.stateCoeff_rvbPEPS}
    \leanok
    \uses{def:peps_rvb_tensor, def:peps_stateCoeff}
    Let $w,h\ge3$. A dimer covering of the torus is a set $c$ of edges such
    that every site lies on exactly one edge of $c$. The RVB PEPS has
    coefficient
    \begin{align}
        \Coeff(\sigma)=\sum_{c}\ \prod_{(v,v+e)\in c}Y_{\sigma_{v+e},\sigma_v},
        \notag
    \end{align}
    the sum running over all dimer coverings, and each edge of $c$ written as
    $(v,v+e)$ with $e\in\{(1,0),(0,1)\}$. The factor $Y_{\sigma_{v+e},\sigma_v}$
    is the coefficient of the singlet $\ket{01}-\ket{10}$ on the two sites, so
    the state is the superposition of all nearest-neighbour singlet coverings
    with the translation-invariant orientation.
\end{theorem}

\begin{proof}\leanok
    \uses{lem:peps_torus_site_tensor_contraction, lem:peps_rvb_local}
    By Lemma~\ref{lem:peps_torus_site_tensor_contraction} the coefficient is a
    sum over bond indices $h_v$, $v_v$ on the right and up edges. By
    Lemma~\ref{lem:peps_rvb_local} a configuration contributes only if the
    edges carrying a non-vacuum index form a dimer covering and the index on
    each such edge is the physical index of the site below or to the left of
    it. Conversely every dimer covering determines one such configuration,
    and its term is the product over sites of the singlet factors on the down
    and left legs; shifting the product by one site gives the formula.
\end{proof}

\section{PEPS from classical models}%
\label{sec:peps_classical_models}

Let $q\ge 1$, let $h\colon\{0,\ldots,q-1\}^{2}\to\mathbb{R}$ be a
nearest-neighbour interaction and $\beta\in\mathbb{R}$ an inverse temperature.

\begin{definition}[Classical energy]%
    \label{def:peps_classical_energy}
    \lean{TNLean.PEPS.classicalEnergy}
    \leanok
    On the torus a configuration $\sigma$ has energy
    \begin{align}
        H(\sigma)=\sum_{v}\bigl(h(\sigma_{v},\sigma_{v-(1,0)})
        +h(\sigma_{v},\sigma_{v-(0,1)})\bigr),
        \notag
    \end{align}
    each nearest-neighbour bond counted once. For symmetric $h$ this is the
    sum of $h$ over unordered bonds.
\end{definition}

\begin{definition}[Classical-model tensor]%
    \label{def:peps_classical_tensor}
    \lean{TNLean.PEPS.classicalBondMatrix, TNLean.PEPS.classicalSiteTensor,
        TNLean.PEPS.classicalPEPS}
    \leanok
    \uses{def:peps_torus_site_tensor}
    Let $M=\sum_{i,j}e^{-\beta h(i,j)/2}|i)(j|$. The classical-model tensor is
    \begin{align}
        A=\Bigl[\sum_{i}\ket{i}(i\,i\,i\,i|\Bigr]
        (\mathbb{1}\otimes\mathbb{1}\otimes M\otimes M):
        \notag
    \end{align}
    the top and right legs copy the physical index $s$, and the down and left
    legs carry $M_{sb}$ and $M_{sl}$ at bond indices $b$ and $l$
    \cite[Appendix~A, ``PEPS from classical models'']{Cirac2021Matrix}. The
    classical-model PEPS places this tensor at every site of the torus.
\end{definition}

\begin{theorem}[Coefficients of the classical-model PEPS]%
    \label{thm:peps_classical_coeff}
    \lean{TNLean.PEPS.stateCoeff_classicalPEPS,
        TNLean.PEPS.conj_mul_stateCoeff_classicalPEPS,
        TNLean.PEPS.normSq_stateCoeff_classicalPEPS}
    \leanok
    \uses{def:peps_classical_tensor, def:peps_classical_energy,
        def:peps_stateCoeff}
    On a torus of width and height at least three the classical-model PEPS
    has coefficient
    $\Coeff(\sigma)=e^{-\beta H(\sigma)/2}$, so
    $|\Coeff(\sigma)|^{2}=e^{-\beta H(\sigma)}$ is the Gibbs weight of
    $\sigma$ \cite{gap:rmp_peps_examples_small_torus}.
\end{theorem}

\begin{proof}\leanok
    \uses{lem:peps_torus_site_tensor_contraction}
    By Lemma~\ref{lem:peps_torus_site_tensor_contraction} the top and right
    legs force every bond index to equal the physical index of the site
    below or to the left of it. The remaining factor at $v$ is
    $M_{\sigma_{v}\sigma_{v-(0,1)}}M_{\sigma_{v}\sigma_{v-(1,0)}}$, one factor
    $e^{-\beta h/2}$ for each bond at $v$ read towards its lower and left
    neighbours, and the product over $v$ is $e^{-\beta H(\sigma)/2}$. The
    coefficient is real, so its squared modulus is $e^{-\beta H(\sigma)}$.
\end{proof}

No bond carries a product of two copies of $M$: the entry
$e^{-\beta h/2}$ is not a square root of a transfer matrix, and no symmetry
of $h$ is used.

\begin{theorem}[Diagonal expectation values are Gibbs averages]%
    \label{thm:peps_classical_expectation}
    \lean{TNLean.PEPS.expectation_classicalPEPS}
    \leanok
    \uses{def:peps_classical_tensor, def:peps_classical_energy,
        def:peps_stateCoeff}
    Let the torus have width and height at least three and let
    $O=\sum_{\sigma}O(\sigma)\ket{\sigma}\bra{\sigma}$ be a diagonal observable.
    For the classical-model PEPS state
    $\ket{\psi}=\sum_{\sigma}\Coeff(\sigma)\ket{\sigma}$,
    \begin{align}
        \frac{\sum_{\sigma}\overline{\Coeff(\sigma)}\,O(\sigma)\Coeff(\sigma)}
        {\sum_{\sigma}\overline{\Coeff(\sigma)}\Coeff(\sigma)}
        =\frac{\sum_{\sigma}O(\sigma)e^{-\beta H(\sigma)}}
        {\sum_{\sigma}e^{-\beta H(\sigma)}},
        \notag
    \end{align}
    that is, $\bra{\psi}O\ket{\psi}/\braket{\psi}{\psi}$ is the expectation
    value of $O$ in the Gibbs state $e^{-\beta H}/Z$
    \cite[Appendix~A, ``PEPS from classical models'']{Cirac2021Matrix}.
\end{theorem}

\begin{proof}\leanok
    \uses{thm:peps_classical_coeff}
    Replace $\overline{\Coeff(\sigma)}\Coeff(\sigma)$ by
    $e^{-\beta H(\sigma)}$ in both sums, by
    Theorem~\ref{thm:peps_classical_coeff}.
\end{proof}

In particular, classical correlation functions are expectation values of
the PEPS. The review adds that at the critical inverse temperature
$\beta_{c}$ the PEPS has correlations decaying as a power law, so its parent
Hamiltonian must be gapless
\cite[Appendix~A, ``PEPS from classical models'']{Cirac2021Matrix}, since a
gapped local Hamiltonian has exponentially decaying correlations
\cite{Hastings2006Spectral,Nachtergaele2006LiebRobinson}.
